\newpage
\hypertarget{para:pirmeu}{}

\mypoemtitle{Pir meu cori alligrari - Stefano Protonotaro}{Sonetto}

\textbf{Commento introduttivo:}\\
Questa canzone di Stefano Protonotaro ci permette di vedere la veste originale degli scritti Siciliani, rappresenta infatti un siciliano illustre con al suo interno elementi di diverse lingue letterarie del tempo tra cui latino, provenzale e siciliano, è per questo motivo abbastanza complicata da comprendere dal punto di vista stilistico

\vspace{1.5em}

\begin{adjustwidth}{2em}{0pt}
\poemlines{5}
\verselinenumbersleft
\begin{verse}
Pir meu cori allegrari, \\
chi multu longiamenti\\
senza alligranza e joi d'amuri è statu,\\
mi ritornu in cantari,\\
ca forsi levimenti\\
da dimuranza turniria in usatu\\
di lu troppu taciri;\\
e quandu l'omu à rasuni di diri,\\
ben di' cantari e mustrari alligranza,\\
ca senza dimustranza\\
joi siria sempri di pocu valuri;\\
dunca ben di' cantari onni amaduri.\\
E si per ben amari\\
cantau juiusamenti\\
homo chi avissi in alcun tempu amatu,\\
ben lu diviria fari\\
plui dilittusamenti\\
eu, chi su di tal donna inamuratu,\\
dundi è dulci placiri,\\
preiu e valenza e juiusu pariri\\
e di billizi cuta[n]t' abondanza,\\
chi illu m'è pir simblanza\\
quandu eu la guardu, sintir la dulzuri\\
chi fa la tigra in illu miraturi;\\
chi si vidi livari\\
multu crudilimenti\\
sua nuritura, chi illa à nutricatu,\\
e si bonu li pari\\
mirarsi dulcimenti\\
dintru unu speclu chi li esti amustratu,\\
chi l'ublia siguiri.\\
Cusì m'è dulci mia donna vidiri:\\
chi 'n lei guardandu met[t]u in ublianza\\
tutta'altra mia intindanza,\\
sì chi instanti mi feri sou amuri\\
d'un culpu chi inavanza tutisuri.\\
Di chi eu putia sanar;\\
multu legeramenti,\\
sulu chi fussi a la mia donna a gratu\\
meu sirviri e pinari;\\
m'eu duitu fortimenti\\
chi quandu si rimembra di sou statu\\
nu lli dia displaciri.\\
Ma si quistu putissi adiviniri,\\
ch'Amori la ferissi de la lanza\\
chi mi fer' e mi lanza,\\
ben crederia guarir de mei doluri,\\
ca sintiramu equalimenti arduri.\\
Purriami laudari\\
d'Amori bonamenti,\\
com'omu da lui beni ammiritatu;\\
ma beni è da blasmari\\
Amur virasementi,\\
quandu illu dà favur da l'unu latu,\\
chi si l'amanti nun sa suffiriri,\\
disia d'amari e perdi sua speranza.\\
Ma eu suf[f]ru in usanza,\\
chi ò vistu adessa bon sufrituri\\
vinciri prova et aquistari hunuri.\\
E si pir suffiriri,\\
ni per amar lialmenti e timiri,\\
homo aquistau d'Amur gran beninanza,\\
digu avir confurtanza\\
eu, chi amu e timu e servi[i] a tutt'uri\\
cilatamenti plu chi altru amaduri.\\
\end{verse}
\end{adjustwidth}


\vspace{0.5em}
{\color{Tomato!40}\hrule height 0.5pt}
\vspace{1em}

\subsubsection*{Analisi}
\begin{quote}
\small 
Nonostante la difficoltà linguistica è interessante mettere in evidenza alcuni elementi. In primo luogo la canzone parla ovviamente della gioia prodotta dall'amore.

Al vv.23-24 cita la tigre vista nei bestiari come un animale vanitoso che si perde guardandosi allo specchio, tanche che nelle battute di caccia i cacciatori mettevano sul terreno degli specchi per far stare fermi gli animali.
Ai vv.35-36 si dice invece che la ferita d'amore è brutta solamente se è uno solo a ferirsi e non tutti e due, significa che se l'amore non è corrisposto fa male ma è anche una metafora alla lancia di Achille che, secondo la leggenda, colpendo una volta ferisce e colpendo nuovamente risana
\end{quote}